% ===========================================================================
%  Préambule — rapport de conduite de projet OC_P13
%
%  Compilation : pdfLaTeX (le compilateur par défaut d'Overleaf).
%  Tous les paquets utilisés sont dans la distribution TeX Live standard,
%  donc rien à installer.
% ===========================================================================

\usepackage[T1]{fontenc}
\usepackage[utf8]{inputenc}
\usepackage[french]{babel}
\usepackage{newtxtext,newtxmath}
\usepackage{microtype}

\usepackage[a4paper,margin=2.4cm,headheight=15pt]{geometry}
\usepackage{graphicx}
\usepackage{booktabs}
\usepackage{tabularx}
\usepackage{array}
\usepackage{longtable}
\usepackage{xcolor}
\usepackage{enumitem}
\usepackage{fancyhdr}
\usepackage{caption}
\usepackage{float}
\usepackage{csquotes}

% Nombres : separateur de milliers par espace fine, virgule decimale.
\usepackage{siunitx}
\sisetup{
  group-separator = {\,},
  group-minimum-digits = 4,
  output-decimal-marker = {,},
  detect-all,
}

\usepackage[hidelinks]{hyperref}
\hypersetup{
  pdftitle={Mise en production d'un modèle de prédiction énergétique},
  pdfauthor={Kevin Lebayle},
  colorlinks=true,
  linkcolor=accent,
  urlcolor=accent,
  citecolor=accent,
}

% --- Couleurs -------------------------------------------------------------
% Les deux cibles du modèle gardent dans le rapport les couleurs qu'elles ont
% dans l'application : bleu pour l'énergie, orange pour les émissions.
\definecolor{accent}{HTML}{2A78D6}
\definecolor{emissions}{HTML}{EB6834}
\definecolor{encadre}{HTML}{F4F6F8}
\definecolor{filet}{HTML}{D5DBE1}
\definecolor{gristexte}{HTML}{52514E}

% --- En-têtes -------------------------------------------------------------
\pagestyle{fancy}
\fancyhf{}
\renewcommand{\headrulewidth}{0.4pt}
\fancyhead[L]{\small\color{gristexte}Conduite de projet — OC\_P13}
\fancyhead[R]{\small\color{gristexte}\leftmark}
\fancyfoot[C]{\small\thepage}

% --- Titres de section ----------------------------------------------------
\usepackage{tikz}
\usetikzlibrary{positioning}

\usepackage{titlesec}
\titleformat{\section}
  {\normalfont\Large\bfseries\color{accent}}{\thesection}{0.8em}{}
\titleformat{\subsection}
  {\normalfont\large\bfseries}{\thesubsection}{0.7em}{}
\titleformat{\subsubsection}
  {\normalfont\normalsize\bfseries\color{gristexte}}{\thesubsubsection}{0.6em}{}

% --- Tableaux -------------------------------------------------------------
\renewcommand{\arraystretch}{1.25}
\newcolumntype{L}[1]{>{\raggedright\arraybackslash}p{#1}}
\newcolumntype{C}[1]{>{\centering\arraybackslash}p{#1}}
\newcolumntype{R}[1]{>{\raggedleft\arraybackslash}p{#1}}

\captionsetup{font=small,labelfont=bf,labelsep=period}

% --- Encadré « à retenir » ------------------------------------------------
% Bâti sur tcolorbox et non sur tabularx : tabularx doit scanner son corps
% jusqu'au \end correspondant, ce qui est impossible quand le \begin est placé
% dans la partie ouvrante d'un \newenvironment. tcolorbox gère ce cas, et
% l'option `breakable` autorise en prime la coupure de page.
\usepackage[most]{tcolorbox}

\newenvironment{constat}[1]
  {\begin{tcolorbox}[
     enhanced, breakable,
     colback=encadre, colframe=encadre,
     borderline west={2.5pt}{0pt}{accent},
     boxrule=0pt, arc=0pt, boxsep=0pt,
     left=12pt, right=10pt, top=8pt, bottom=8pt,
   ]%
   \textbf{#1}\par\vspace{0.35em}}
  {\end{tcolorbox}}

% --- Raccourcis -----------------------------------------------------------
\newcommand{\code}[1]{\texttt{\small #1}}
\newcommand{\ecart}[1]{\textcolor{accent}{\textbf{#1}}}
